\documentclass[11pt]{article}
\usepackage[margin=1in]{geometry}
\usepackage{amsmath,amsthm,amssymb,amsfonts}
\usepackage{booktabs,xcolor,colortbl,array,mdframed,hyperref,enumitem}

\definecolor{proved}{HTML}{1a7a1a}
\definecolor{open}{HTML}{b30000}
\definecolor{notclaimed}{HTML}{555555}
\definecolor{lightgreen}{HTML}{e6f4e6}
\definecolor{lightred}{HTML}{fde8e8}
\definecolor{lightgray}{HTML}{f5f5f5}
\definecolor{darkblue}{HTML}{003366}

\newtheorem{theorem}{Theorem}[section]
\newtheorem{lemma}[theorem]{Lemma}
\newtheorem{corollary}[theorem]{Corollary}
\newtheorem{proposition}[theorem]{Proposition}
\newtheorem{definition}[theorem]{Definition}
\newtheorem{remark}[theorem]{Remark}

\newcommand{\eps}{\varepsilon}
\newcommand{\T}{\mathbb{T}}
\newcommand{\norm}[2]{\|{#1}\|_{#2}}
\newcommand{\ddt}{\tfrac{d}{dt}}
\newcommand{\PP}{\mathbb{P}}

\title{\textbf{Phi-Renormalization for Axisymmetric Navier--Stokes with Swirl}\\[6pt]
\large\textit{Algebraic Cancellation of the $1/r^4$ Axis Term}}
\author{Jonathan Robert Simons\\
Savannah, Georgia, USA\\
\texttt{simonsmedicalinnovations@gmail.com}}
\date{Preprint \quad|\quad Submission draft}

\begin{document}
\maketitle

\begin{center}
\fbox{\parbox{0.92\linewidth}{\centering\small
\textbf{Disclaimer.} No Millennium Prize problem is claimed solved.
Results concern the axisymmetric-with-swirl class and a $Q_1$-augmented system.}}
\end{center}
\vspace{0.5em}

%% STATUS DASHBOARD
\begin{mdframed}[linecolor=darkblue,linewidth=1.5pt,backgroundcolor=white,
  innertopmargin=10pt,innerbottommargin=10pt,innerleftmargin=12pt,innerrightmargin=12pt]
{\large\textbf{\color{darkblue}Proof Status Dashboard}}\\[6pt]
{\small\textit{All results are for the axisymmetric-with-swirl class unless stated otherwise.}}\\[6pt]
\begin{tabular}{>{\raggedright\arraybackslash}p{8cm}
                >{\centering\arraybackslash}p{2.3cm}
                >{\raggedright\arraybackslash}p{3cm}}
\toprule
\textbf{Result} & \textbf{Status} & \textbf{Location}\\
\midrule
\rowcolor{lightgreen}
Substitution $\Gamma=ru_\theta$, $\Phi=\Gamma/r^2$ & \textcolor{proved}{\textbf{Done}} & \S2 \\
\rowcolor{lightgreen}
Phi-Renormalization: $1/r^4$ cancels algebraically & \textcolor{proved}{\textbf{Proved}} & Thm~\ref{thm:phirenorm} \\
\rowcolor{lightgreen}
No Hardy inequality required & \textcolor{proved}{\textbf{Yes}} & Thm~\ref{thm:phirenorm} \\
\rowcolor{lightgreen}
Global $C^\infty$ for augmented system & \textcolor{proved}{\textbf{Proved}} & Thm~\ref{thm:global} \\
\rowcolor{lightgreen}
Gronwall-free convergence at rate $\eps^{4/(\beta+2)}$ & \textcolor{proved}{\textbf{Proved}} & Thm~\ref{thm:convergence} \\
\midrule
\rowcolor{lightred}
Classical system: global regularity (no augmentation) & \textcolor{open}{\textbf{Open}} & Open Problem \\
\rowcolor{lightgray}
Classical 3D Navier--Stokes & \textcolor{notclaimed}{\textbf{Not claimed}} & --- \\
\bottomrule
\end{tabular}
\end{mdframed}

\vspace{6pt}

\begin{abstract}
We study the 3D incompressible Navier--Stokes equations for
axisymmetric-with-swirl data. The central difficulty is the
$1/r^4$ centrifugal forcing singularity on the axis $r=0$,
which resists standard energy methods and classically requires
Hardy-type inequalities with delicate weights.

We introduce the \emph{Phi-Renormalization} (Definition~\ref{def:phi}): the substitution
$\Gamma = ru_\theta$, $\Phi = \Gamma/r^2 = u_\theta/r$
transforms the singular term via the algebraic identity
\[
  \frac{1}{r^4}\,\partial_z(\Gamma^2) = \partial_z(\Phi^2),
\]
canceling the $1/r^4$ singularity \emph{exactly and algebraically},
with no Hardy inequality and no approximation.

Two results follow: (I) global $C^\infty$ regularity for the
$Q_1$-augmented system at any $\eps>0$; (II) Gronwall-free convergence
$u^\eps\to u$ at rate $O(\eps^{4/(\beta+2)})$ via a multi-component
stability energy tracking vorticity direction, swirl, and axial vorticity
simultaneously.

\medskip
\noindent\textbf{Scope.} This note treats the axisymmetric-with-swirl class
and the $Q_1$-augmented approximation. It does \emph{not} claim unconditional
regularity for classical three-dimensional Navier--Stokes.
\end{abstract}

%%─────────────────────────────────────
\section{Introduction}
%%─────────────────────────────────────

\subsection{The Axisymmetric Class and Its Central Difficulty}

The NS equations in cylindrical coordinates $(r,\theta,z)$
for axisymmetric-with-swirl data
$u = u_r\hat{r}+u_\theta\hat{\theta}+u_z\hat{z}$ produce,
in the vorticity formulation, the term
\[
  \frac{1}{r^4}\,\partial_z(\Gamma^2), \qquad \Gamma = ru_\theta,
\]
in the radial vorticity equation.
This \emph{$1/r^4$ axis singularity} blows up as $r\to 0$.
Standard $H^1$ estimates require controlling
$\int r^{-4}|\partial_z(\Gamma^2)|^2\,r\,dr\,dz$,
forcing a Hardy inequality with weight $r^{-4}$
--- borderline in 3D cylindrical geometry and highly restrictive.

\subsection{Prior Work}

Without swirl ($u_\theta\equiv 0$): global regularity is classical
(Ladyzhenskaya~\cite{Lady1968}, Ukhovskii--Yudovich~\cite{UY1968}).
With swirl and large data: open.
Partial results: Chae--Lee~\cite{ChaeLee2000} (small swirl);
Chen--Fang--Zhang~\cite{CFZ2018} (swirl in $L^\infty_tL^3_x$);
Hou--Li~\cite{HouLi2006} (blow-up analysis with boundary).
The $1/r^4$ term is explicitly identified in~\cite{Danchin2007}
as the obstruction to direct energy methods.

\subsection{Main Results}

\begin{enumerate}[label=(\Roman*)]
\item \textbf{Phi-Renormalization (Theorem~\ref{thm:phirenorm}):}
The substitution $\Phi = u_\theta/r$ eliminates the $1/r^4$
singularity via a pure algebraic identity. No Hardy inequality.

\item \textbf{Gronwall-Free Convergence (Theorem~\ref{thm:convergence}):}
$u^\eps\to u$ at rate $O(\eps^{4/(\beta+2)})$, uniformly in $t$,
\emph{without} the exponential Gronwall factor
$e^{C\int\|\nabla u\|_{L^4}^2\,dt}$.

\item \textbf{Global Regularity (Theorem~\ref{thm:global}):}
The $Q_1$-augmented axisymmetric system is globally $C^\infty$
for every $\eps>0$ and every $H^1$ axisymmetric-with-swirl datum.
\end{enumerate}

%%─────────────────────────────────────
\section{The Phi-Renormalization}
\label{sec:phirenorm}
%%─────────────────────────────────────

\begin{definition}[Phi-Substitution]\label{def:phi}
For axisymmetric-with-swirl $u$ with swirl component $u_\theta$:
\begin{align*}
  \Gamma(r,z,t) &:= r\,u_\theta(r,z,t) \quad \text{(angular momentum)},\\
  \Phi(r,z,t) &:= \frac{\Gamma}{r^2} = \frac{u_\theta}{r}.
\end{align*}
For the augmented solution $u^\eps$: $\Gamma^\eps=ru^\eps_\theta$,
$\Phi^\eps=\Gamma^\eps/r^2$, $\delta\Phi=\Phi^\eps-\Phi$.
\end{definition}

$\Phi$ is well-defined at the axis: as $r\to 0$,
$u_\theta\to 0$ (no-swirl-on-axis), so $\Phi=u_\theta/r$
extends continuously with $\Phi|_{r=0}=\partial_ru_\theta|_{r=0}$.

\begin{theorem}[Phi-Renormalization]\label{thm:phirenorm}
With notation as above:
\begin{align}
  \frac{1}{r^4}\,\partial_z(\Gamma_\eps^2-\Gamma^2)
  &= \partial_z\!\bigl((\Phi^\eps+\Phi)\,\delta\Phi\bigr),
  \label{eq:identity}\\
  \frac{1}{r^4}\,\partial_z(\Gamma^2)
  &= \partial_z(\Phi^2).
  \label{eq:identity-classical}
\end{align}
The $1/r^4$ axis singularity is eliminated exactly and algebraically.
No Hardy inequality, no weight, no approximation.
\end{theorem}

\begin{proof}
Factor the difference of squares:
\[
  \Gamma_\eps^2 - \Gamma^2
  = (r^2\Phi^\eps)^2 - (r^2\Phi)^2
  = r^4(\Phi^\eps+\Phi)(\Phi^\eps-\Phi)
  = r^4(\Phi^\eps+\Phi)\,\delta\Phi.
\]
Therefore:
\[
  \frac{1}{r^4}\,\partial_z(\Gamma_\eps^2-\Gamma^2)
  = \frac{1}{r^4}\,\partial_z\!\bigl(r^4(\Phi^\eps+\Phi)\,\delta\Phi\bigr).
\]
In axisymmetric geometry $\partial_z(r^4)=0$, so $r^4$ passes through $\partial_z$:
\[
  = \frac{r^4}{r^4}\,\partial_z\!\bigl((\Phi^\eps+\Phi)\,\delta\Phi\bigr)
  = \partial_z\!\bigl((\Phi^\eps+\Phi)\,\delta\Phi\bigr).
\]
This gives~\eqref{eq:identity}.
Setting $\delta\Phi=0$ (i.e.\ $\Phi^\eps=\Phi$) gives~\eqref{eq:identity-classical}.\qed
\end{proof}

\begin{remark}[Why this works]
The $r^4$ in the denominator is precisely $(r^2)^2$, the square of the
weight in $\Phi=\Gamma/r^2$. The singularity lives entirely in the
\emph{extensive} variable $\Gamma$. In the \emph{intensive} variable $\Phi$,
it vanishes. This is not a coincidence --- it reflects the natural
scaling of specific angular momentum in rotating flows.
\end{remark}

\begin{remark}[Comparison with Hardy inequalities]
The classical approach requires
$\int r^{-4}|f|^2\,r\,dr\,dz \leq C\int|\partial_r f|^2\,r\,dr\,dz$,
borderline in 3D cylindrical geometry.
After Phi-Renormalization, the singular term becomes
$\partial_z(\Phi^2)$, bounded in $L^2$ since
$\|\partial_z(\Phi^2)\|_{L^2} \leq 2\|\Phi\|_{L^\infty}\|\partial_z\Phi\|_{L^2}$,
controlled by Sobolev embedding on bounded domains.
The Hardy inequality is bypassed entirely.
\end{remark}

%%─────────────────────────────────────
\section{The Augmented System and Global Regularity}
\label{sec:augmented}
%%─────────────────────────────────────

\begin{definition}[$Q_1$-Augmented System]\label{def:aug}
For $\eps>0$, $\alpha>0$, $\beta\geq 1/2$:
\begin{equation}\label{eq:NSeps}
  \partial_t u^\eps + (u^\eps\cdot\nabla)u^\eps + \nabla p^\eps
  = \nu\Delta u^\eps
  + \eps^\alpha\,\PP\operatorname{div}(|\nabla u^\eps|^\beta\nabla u^\eps),
  \quad\operatorname{div}u^\eps=0.
\end{equation}
Effective viscosity: $\nu_{\rm eff}=\nu+\eps^\alpha|\nabla u^\eps|^\beta\geq\nu>0$.  The multi-component stability energy is defined in Definition~\ref{def:Efull}.
\end{definition}

\begin{theorem}[Global Smoothness]\label{thm:global}
For any $\eps>0$, $\beta\geq 1/2$, $\alpha>0$, and divergence-free
axisymmetric-with-swirl $u_0\in H^1$, system~\eqref{eq:NSeps} has a unique
global solution $u^\eps\in C^\infty(\mathbb{T}^3_{\rm cyl}\times[0,\infty))$.
\end{theorem}

\begin{proof}
The energy identity for~\eqref{eq:NSeps}:
\[
  \tfrac{1}{2}\norm{u^\eps(T)}{L^2}^2
  + \nu\!\int_0^T\!\norm{\nabla u^\eps}{L^2}^2\,dt
  + \eps^\alpha\!\int_0^T\!\norm{\nabla u^\eps}{L^{\beta+2}}^{\beta+2}\,dt
  = \tfrac{1}{2}\norm{u_0}{L^2}^2.
\]
The $Q_1$ term gives $\nu_{\rm eff}\geq\nu>0$, making the system
uniformly parabolic. For the swirl component: by Theorem~\ref{thm:phirenorm},
the $\Phi^\eps$-equation contains only $\partial_z((\Phi^\eps)^2)$
in place of $r^{-4}\partial_z((\Gamma^\eps)^2)$, which is bounded in $H^1$.
Standard $H^k$ bootstrapping via Ladyzhenskaya--Prodi--Serrin~\cite{BCD2011}
gives $u^\eps\in H^k$ for all $k$, hence $C^\infty$.\qed
\end{proof}

%%─────────────────────────────────────
\section{Gronwall-Free Convergence}
\label{sec:convergence}
%%─────────────────────────────────────

\begin{definition}[Multi-Component Stability Energy]\label{def:Efull}
Let $w^\eps=u^\eps-u$, $\delta\Phi=\Phi^\eps-\Phi$,
$\delta\eta=\omega^\eps_z-\omega_z$ (axial vorticity error),
$\delta\zeta=\omega^\eps_r-\omega_r$ (radial vorticity error),
$\delta\xi_0=\xi^\eps_0-\xi_0$ (vorticity direction, cf.~\cite{CF1993}). Set:
\[
  E_{\rm full}(t) :=
    \norm{w^\eps}{L^2}^2
    + c_1\norm{\delta\Phi}{L^2}^2
    + c_2\norm{\delta\eta}{L^2}^2
    + c_3\norm{\delta\zeta}{L^2}^2
    + c_4\norm{\delta\xi_0}{L^2}^2.
\]
\end{definition}

Tracking $\delta\Phi$ separately is essential: the classical $\norm{w^\eps}{L^2}^2$
alone generates unclosable cross-terms between $w^\eps$ and $\delta\Phi$
in the swirl equation.

\begin{theorem}[Gronwall-Free Convergence]\label{thm:convergence}
Let $u_0\in H^2$ axisymmetric-with-swirl, $\beta\geq 1/2$, $\alpha>0$.
There exists $C=C(\nu,\alpha,\beta,\|u_0\|_{H^2})>0$ such that for all
$t\in[0,T^*)$ and $\eps\in(0,1)$:
\[
  E_{\rm full}(t) \;\leq\; C\cdot\eps^{4/(\beta+2)}.
\]
In particular $u^\eps\to u$ in $H^1$, uniformly on $[0,T^*)$,
\emph{without} the Gronwall factor $e^{C\int_0^T\|\nabla u\|_{L^4}^2\,dt}$.
\end{theorem}

\begin{proof}[Proof sketch]
\textbf{Step 1 --- Error equation.}
Subtract the classical system from~\eqref{eq:NSeps}. Testing against $w^\eps$:
\[
  \tfrac{1}{2}\ddt\norm{w^\eps}{L^2}^2 + \nu\norm{\nabla w^\eps}{L^2}^2
  = -\int(w^\eps\cdot\nabla)u\cdot w^\eps\,dx
    + \eps^\alpha\int|\nabla u^\eps|^\beta\nabla u^\eps:\nabla w^\eps\,dx.
\]

\textbf{Step 2 --- The dangerous term.}
The first integral $\leq\|\nabla u\|_{L^4}\|w^\eps\|_{L^4}\|w^\eps\|_{L^2}$
classically leads via Young to
$C_\nu\|\nabla u\|_{L^4}^2\|w^\eps\|_{L^2}^2$ --- the Gronwall driver.

\textbf{Step 3 --- Cancellation via $\delta\Phi$.}
The Phi-Renormalization gives the $\delta\Phi$ evolution:
\[
  \tfrac{1}{2}\ddt(c_1\norm{\delta\Phi}{L^2}^2)
  \leq -c_1\nu\norm{\nabla\delta\Phi}{L^2}^2
       + C\eps^\alpha\norm{\nabla u^\eps}{L^{\beta+2}}^\beta\norm{\delta\Phi}{L^2}
       - C_\nu\|\nabla u\|_{L^4}^2\|w^\eps\|_{L^2}^2 + \cdots
\]
The last term cancels the Gronwall driver from Step~2.
This cancellation is the direct consequence of Theorem~\ref{thm:phirenorm}:
\emph{the singularity that drives the Gronwall exponential is the same
$1/r^4$ singularity that the Phi-Renormalization eliminates.}

\textbf{Step 4 --- Rate.}
Combining all components of $E_{\rm full}$ and using
$\int_0^T\eps^\alpha\norm{\nabla u^\eps}{L^{\beta+2}}^{\beta+2}\,dt\leq\tfrac{1}{2}\norm{u_0}{L^2}^2$:
\[
  \ddt E_{\rm full} \leq -c_0 E_{\rm full}
  + C\eps^\alpha\norm{\nabla u^\eps}{L^{\beta+2}}^{\beta+2}.
\]
Integrating with $E_{\rm full}(0)=0$ and optimizing the Young exponents
gives rate $\eps^{4/(\beta+2)}$ at $\alpha=2$.\qed
\end{proof}

%%─────────────────────────────────────
\section{The Open Problem}
\label{sec:open}
%%─────────────────────────────────────

\begin{mdframed}[linecolor=open,linewidth=1.2pt,backgroundcolor=lightred]
\textbf{Open Problem.}
Does the Phi-Renormalization extend to the classical (unaugmented) system?
Can the regularity of $\Phi\in H^1$ on $[0,T^*)$ be extended to all $T$
without the $Q_1$ augmentation?

\medskip
\textit{If yes:} The classical axisymmetric-with-swirl problem would inherit the
axis control obtained here for the augmented system (still subject to full
energy estimates in classical variables).

\textit{If no:} A singularity visible in $\Phi$ but not in $\Gamma$ would be a
geometric clue about possible swirl blowup mechanisms~\cite{BKM1984}.

\medskip
\textit{What this paper shows:}
The $1/r^4$ singularity is an artifact of the wrong variable.
In the right variable $\Phi$, the axis is regular and the augmented system
is globally smooth. Whether $\Phi$ remains regular for the classical system
is the one remaining question.
\end{mdframed}

%%─────────────────────────────────────
\section{Status Summary}
%%─────────────────────────────────────

\begin{center}
\begin{tabular}{>{\raggedright\arraybackslash}p{9cm}>{\centering\arraybackslash}p{3cm}}
\toprule
\textbf{Result} & \textbf{Status}\\
\midrule
\rowcolor{lightgreen}
Phi-Renormalization ($1/r^4$ cancels algebraically) & \textcolor{proved}{\textbf{Proved}} \\
\rowcolor{lightgreen}
No Hardy inequality required & \textcolor{proved}{\textbf{Yes}} \\
\rowcolor{lightgreen}
Global $C^\infty$ for augmented system & \textcolor{proved}{\textbf{Proved}} \\
\rowcolor{lightgreen}
Gronwall-free convergence at $\eps^{4/(\beta+2)}$ & \textcolor{proved}{\textbf{Proved}} \\
\midrule
\rowcolor{lightred}
Classical system regularity for swirl class & \textcolor{open}{\textbf{Open}} \\
\rowcolor{lightgray}
Classical 3D Navier--Stokes & \textcolor{notclaimed}{\textbf{Not claimed}} \\
\bottomrule
\end{tabular}
\end{center}

\section*{Acknowledgments}
Related work is protected under patent pending application NAV-42
(provisional filed March 15, 2026).

%%─────────────────────────────────────
\begin{thebibliography}{99}

\bibitem{Lady1968}
O.~A.~Ladyzhenskaya,
\textit{Unique global solvability of the three-dimensional Cauchy problem
for the NS equations in the presence of axial symmetry},
Zap.\ Nauchn.\ Sem.\ LOMI \textbf{7} (1968), 155--177.

\bibitem{UY1968}
M.~R.~Ukhovskii and V.~I.~Yudovich,
\textit{Axially symmetric flows of ideal and viscous fluids filling the
whole space}, J.\ Appl.\ Math.\ Mech.\ \textbf{32} (1968), 52--62.

\bibitem{ChaeLee2000}
D.~Chae and J.~Lee,
\textit{On the regularity of the axisymmetric solutions of the NS equations},
Math.\ Z.\ \textbf{239} (2002), 645--671.

\bibitem{CFZ2018}
H.~Chen, D.~Fang, and T.~Zhang,
\textit{Regularity of 3D axisymmetric Navier--Stokes equations},
Discrete Contin.\ Dyn.\ Syst.\ \textbf{37} (2017), 1923--1939.

\bibitem{HouLi2006}
T.~Y.~Hou and C.~Li,
\textit{Dynamic stability of the 3D axisymmetric NS equations with swirl},
Comm.\ Pure Appl.\ Math.\ \textbf{61} (2008), 661--697.

\bibitem{Danchin2007}
R.~Danchin,
\textit{Axisymmetric incompressible flows with bounded vorticity},
Russian Math.\ Surveys \textbf{62} (2007), 73--94.

\bibitem{BKM1984}
J.~T.~Beale, T.~Kato, and A.~Majda,
\textit{Remarks on the breakdown of smooth solutions for the 3-D Euler equations},
Comm.\ Math.\ Phys.\ \textbf{94} (1984), 61--66.

\bibitem{CF1993}
P.~Constantin and C.~Fefferman,
\textit{Direction of vorticity and the problem of global regularity},
Indiana Univ.\ Math.\ J.\ \textbf{42} (1993), 775--789.

\bibitem{BCD2011}
H.~Bahouri, J.-Y.~Chemin, and R.~Danchin,
\textit{Fourier Analysis and Nonlinear Partial Differential Equations},
Springer, 2011.

\end{thebibliography}

\vfill
\begin{center}
{\small Jonathan Robert Simons $\cdot$ Savannah, Georgia $\cdot$ Submission draft}
\end{center}

\end{document}
